\chapter{Urinalysis: Ketones}
\section*{What Is This Test?}
Urine ketones testing detects the presence of ketone bodies (acetoacetate, beta-hydroxybutyrate, and acetone) in the urine. Ketone bodies are produced by hepatic mitochondrial fatty acid beta-oxidation and ketogenesis when insulin is deficient or absent and counter-regulatory hormones (glucagon, cortisol, epinephrine, growth hormone) are elevated, as occurs in: diabetic ketoacidosis (DKA), prolonged fasting/starvation, alcoholic ketoacidosis, very low-carbohydrate ketogenic diets, and pregnancy (starvation ketosis). In DKA, insulin deficiency results in uncontrolled lipolysis in adipose tissue, releasing free fatty acids that are taken up by the liver and converted to ketone bodies. The accumulation of acetoacetate and beta-hydroxybutyrate produces an increased anion gap metabolic acidosis. The urine ketone dipstick uses the Legal test (nitroprusside reaction): sodium nitroprusside reacts with acetoacetate and acetone (but NOT beta-hydroxybutyrate) at an alkaline pH to produce a purple color. The dipstick detects acetoacetate with a sensitivity threshold of 5--10 mg/dL. In DKA, the predominant ketone body is beta-hydroxybutyrate (which is NOT detected by the dipstick). As DKA resolves with insulin therapy, beta-hydroxybutyrate is oxidized to acetoacetate, paradoxically causing the urine ketone test to remain positive or even increase (``ketone paradox'') despite clinical improvement. For this reason, direct measurement of serum beta-hydroxybutyrate is preferred for monitoring DKA treatment.
\section*{Alternative Names}
Urine Ketones; Ketonuria; Ketone Bodies in Urine; Acetoacetate Test; Legal's Test
\section*{Principle}
Nitroprusside reaction (Legal test): Sodium nitroprusside + acetoacetate (or acetone) ---> purple-colored complex in alkaline pH. The dipstick does NOT react with beta-hydroxybutyrate, the predominant ketone body in DKA (accounting for 75--80\% of total ketone bodies in severe DKA). The detection threshold is 5--10 mg/dL. False-positive: highly pigmented urine, levodopa metabolites, sulfhydryl compounds including captopril and mesna (2-mercaptoethane sulfonate). False-negative: improperly stored urine specimens in which acetoacetate has decarboxylated to acetone and evaporated.
\section*{History}
Acetone was first detected in the urine of a diabetic patient by Wilhelm Petters in 1857, and in 1883 Nicholas Legal described the sodium nitroprusside reaction for acetoacetate that still underlies the modern dipstick. In the early 1900s, Magnus-Levy established the relationship between ketone bodies and diabetic acidosis. Following the introduction of insulin in 1922, ketone measurement became essential for recognizing and monitoring diabetic ketoacidosis. Helen Free and Alfred Free developed the first practical urine ketone dipstick (Ketostix) in the 1950s, making the nitroprusside test widely available at the bedside. Recognition that the dipstick does not detect beta-hydroxybutyrate --- the predominant ketone in DKA --- led to the development of serum beta-hydroxybutyrate meters in the 1990s for more accurate monitoring of ketosis.
\section*{Why This Test Is Ordered}
\begin{itemize}
  \item Diagnosis of diabetic ketoacidosis in patients with hyperglycemia, and differentiation from hyperosmolar hyperglycemic state
  \item Monitoring of DKA treatment, where a fall in serum beta-hydroxybutyrate confirms resolution even while urine ketones lag behind (ketone paradox)
  \item Evaluation of metabolic acidosis with an increased anion gap to identify a ketotic component
  \item Assessment of ketosis in prolonged fasting, very low-carbohydrate (ketogenic) diets, and pregnancy (starvation ketosis)
  \item Diagnosis of alcoholic ketoacidosis in patients with chronic alcohol use, vomiting, and starvation
  \item Monitoring of patients on SGLT2 inhibitors presenting with suspected euglycemic ketoacidosis
\end{itemize}
\section*{Is This a Primary or Secondary Test}
Urine ketone testing is a primary point-of-care test for detecting ketosis in DKA and is a rapid, inexpensive screen. Because it does not detect beta-hydroxybutyrate, serum beta-hydroxybutyrate is the preferred secondary (confirmatory and monitoring) test during treatment. The urine dipstick remains useful for screening, but a negative urine ketone test does not exclude DKA when beta-hydroxybutyrate predominates.
\section*{How This Test Is Performed}
A freshly voided urine sample is collected in a clean container, and the reagent strip is dipped into the specimen or passed through the midstream urine. The reagent pad is compared against the color chart after the manufacturer-specified time (typically 15--60 seconds). Results are reported semiquantitatively as negative, small, moderate, or large (approximate mg/dL of acetoacetate). For DKA monitoring, a capillary blood sample is tested with a beta-hydroxybutyrate meter.
\section*{What Preparation Is Needed}
None. A fresh, well-mixed urine sample is preferred because acetoacetate is unstable; standing at room temperature causes decarboxylation to acetone, which volatilizes and produces false-negative results. Testing should be performed within 1--2 hours of collection.
\section*{Contraindications and Precautions}
No contraindications to urine collection. Precautions: (1) false-positive results occur with highly pigmented urine, levodopa metabolites, and sulfhydryl compounds (captopril, mesna); (2) false-negative results occur with old or improperly stored specimens; (3) highly acidic or strongly dilute urine can blunt the reaction; (4) the dipstick does not detect beta-hydroxybutyrate, so a negative urine test does not exclude DKA; (5) positive results should be interpreted in clinical context, because some drugs cross-react with the reagent pad.
\section*{Risks}
No physical risk; the test is noninvasive. The main concern is misinterpretation when the specimen is stale or when the beta-hydroxybutyrate-predominant state of severe DKA is not recognized.
\section*{What Happens After the Test}
A positive urine ketone test in a diabetic patient triggers measurement of serum glucose, beta-hydroxybutyrate, electrolytes (anion gap, bicarbonate), and blood gas to confirm DKA. During treatment, serum beta-hydroxybutyrate is rechecked until it falls below 1.0--1.5 mmol/L and the anion gap normalizes. Urine ketones may remain positive for 12--24 hours after clinical resolution (ketone paradox) and should not delay transition from intravenous insulin when other parameters have improved.
\section*{Understanding Results}
\subsection*{Negative}
No ketones detected. Normal metabolism without significant fatty acid oxidation or ketogenesis. Well-controlled diabetes.
\subsection*{Positive}
Small (5--15 mg/dL): Fasting, acute illness with decreased oral intake, vomiting, ketogenic diet, early DKA. Moderate (15--40 mg/dL): Early DKA, alcoholic ketoacidosis, pregnancy with inadequate caloric intake. Large (40--80--160+ mg/dL): DKA, severe alcoholic ketoacidosis. In DKA: positive urine ketones confirm the diagnosis, but serum beta-hydroxybutyrate is the preferred biomarker for quantifying ketosis severity and monitoring treatment response (goal: beta-hydroxybutyrate <1.0 mmol/L before transition from IV insulin infusion). The urine ketone test may remain positive for 12--24 hours after resolution of DKA due to the ketone paradox.
\section*{Normal Results}
Negative (urine acetoacetate <5--10 mg/dL). Normal individuals do not excrete significant quantities of ketones. Pregnant women may have mild ketonuria with prolonged fasting or vomiting (hyperemesis gravidarum).
\section*{Follow-up Tests}
Serum beta-hydroxybutyrate (<0.6 mmol/L normal, >1.0 mmol/L = hyperketonemia, >3.0 mmol/L = severe ketosis/DKA), serum glucose, basic metabolic panel (anion gap, bicarbonate), serum osmolality, arterial or venous blood gas (pH, PCO2, HCO3), and serum lactate.
\section*{References}
\begin{enumerate}
  \item Kitabchi AE, Umpierrez GE, Murphy MB, et al. Management of hyperglycemic crises in patients with diabetes. Diabetes Care. 2001;24(1):131--153.
  \item Laffel L. Ketone bodies: a review of physiology, pathophysiology and application to monitoring in diabetes. Diabetes Metab Res Rev. 1999;15(6):412--426.
  \item Kitabchi AE, Umpierrez GE, Miles JM, Fisher JN. Hyperglycemic crises in adult patients with diabetes. Diabetes Care. 2009;32(7):1335--1343.
  \item McGarry JD, Foster DW. Regulation of hepatic fatty acid oxidation and ketone body production. Annu Rev Biochem. 1980;49:395--420.
  \item Burtis CA, Ashwood ER, Bruns DE. Tietz Textbook of Clinical Chemistry and Molecular Diagnostics. 5th ed. Elsevier; 2012.
\end{enumerate}
\clearpage
